\section{A shared index for maximum cyclic relator shortening}
\label{sec:joint-cyclic-overlap}

\subsection{The bottleneck and the exact query}

The existing native positive-certificate path reconstructs the complete
Wirtinger presentation, performs generator eliminations and Whitehead moves,
and searches for a length-reducing relator substitution when these moves
stall.  The relevant local query takes an indexed list
$(R_1,\ldots,R_s)$ of explicit cyclic words.  Let $m$ be the number of
nonempty slots and $L=\sum_i |R_i|$.  For distinct slots $i\ne j$, rotate
$R_i$ and either $R_j$ or $R_j^{-1}$ to write
\[
 R_i\sim UW,\qquad R_j^{\varepsilon}\sim UV,
 \qquad \varepsilon\in\{1,-1\}.
\]
Replacing the target by $V^{-1}W$ while retaining the donor preserves the
normal closure.  Indeed, $(UV)^{-1}(UW)=V^{-1}W$ and
$(UV)(V^{-1}W)=UW$.  The guaranteed decrease before incidental free
cancellation is
\[
 g=2|U|-|R_j|.
\]
The search must find the maximum positive $g$ over all target rotations,
all donor rotations and both donor signs, or prove that none exists.
No small-cancellation hypothesis is needed for the soundness of a move.
Stalling is inconclusive for unknot recognition.

The maintained pairwise implementation builds a suffix automaton for each
oriented doubled donor and scans every eligible doubled target.  Its bound
is $O(s+mL)$ dictionary operations and $O(L)$ auxiliary space.  Profiling
the actual 141-crossing Gordian input found one such query after 129 generator
eliminations.  Twelve nonempty relators then have total length $17\,504$.
This query, rather than the elimination loop, was the dominant producer
cost.  The new query removes the factor $m$ from its asymptotic operation
bound; it does not bound the number of subsequent presentation moves.

\subsection{Cyclic windows as capped points on a suffix-link tree}

A nonempty word $w$ of length $d$ has all its cyclic rotations among the
length-$d$ windows ending at positions
\[
 d-1,d,\ldots,2d-2
\]
of $ww[0:d-1]$, using zero-based positions.  There are exactly $d$ such
windows, even when several spell the same word.  Each shorter cyclic
segment is a suffix of one of them.  A capacity $d$ prevents any accepted
match from wrapping around the donor or target twice.

Build one suffix automaton for the concatenation of the positive doubled
relators, with a separator after each block.  In the implementation the
separator is a fresh object outside the signed-generator alphabet.  The
indexed string has length exactly $2L$.  Standard suffix-automaton theory
provides $O(L)$ states and transitions; the classical linear-size subword
index is due to Blumer and coauthors.  A state has a maximum represented
substring length $\lambda(v)$ and a suffix link $p(v)$ satisfying
$\lambda(p(v))<\lambda(v)$.  The represented lengths on its suffix-link
edge are the integers
\[
 \lambda(p(v))+1,\ldots,\lambda(v).
\]
All strings represented by one state are suffixes of its unique longest
string.  Consequently a virtual point at an intermediate edge depth has
a well-defined word, not merely a numerical length.  Thus an occurrence
of a word of length $d$ can be represented by a point
at depth $d$ on the suffix-link path of its ending state.  Insert that
point as a terminal cap if it is not already a state vertex.

For every positive rotation, mark its capped point with two roles: it is
a possible target and a possible positive donor.  The mark retains its
original relator slot, its word length, its sign, and the ending position
inside the doubled source.  Equal words in distinct slots retain distinct
identities.  Losing this information would permit an invalid self-donor
substitution.

The inverse donors need not enlarge the index.  Scan each inverse doubled
relator through the already completed automaton, maintaining the longest
matched suffix.  At each of its $d$ cyclic-window endpoints, let $h$ be
the length of that matched suffix.  If $h>0$, attach an inverse-donor mark
at depth $\min(d,h)$, but retain the original donor length $d$ for the gain
calculation.  This is a partially represented inverse window, not a claim
that its full length-$d$ rotation occurs in the indexed text.

\begin{lemma}[Completeness of the inverse-window marks]
Every common cyclic segment of an inverse donor and a positive target is
represented by the common ancestral path of one inverse-donor mark and one
positive-target mark.  Conversely, every such common path denotes a literal
common segment of those oriented relators.
\end{lemma}
\begin{proof}
Let the segment have length $u$.  Choose the donor's length-$d$ cyclic
window ending with this segment; its endpoint is one of the $d$ recorded
endpoints.  Since the segment occurs in a positive target, it occurs in
the indexed text.  The longest suffix match at that donor endpoint has
length at least $u$, and $u\le d$, so the inverse mark has depth at least
$u$.  Its suffix-link path therefore represents the required segment.
The positive target has an analogous complete cyclic-window mark.
Conversely, an ancestor of a mark represents a suffix of the actual marked
window.  The depth caps ensure that this suffix has length at most the
corresponding original relator length.  A common ancestor represents the
same substring in both windows.
\end{proof}

The capped tree contains at most $2L$ marked windows and at most $2L$ new
cap vertices.  Together with the automaton states it has $O(L)$ vertices.
Marks can coincide, in which case their original slots remain separate
records.  No quadratic list of explicit rotations is constructed.

\subsection{Two representatives suffice for the distinct-slot constraint}

At a capped-tree vertex of depth $h$, all descendant marks contain the same
length-$h$ suffix.  A donor of length $d$ and any distinct-slot target
therefore supply gain $2h-d$.  At this fixed depth, maximizing gain means
minimizing donor length among eligible donor slots.

\begin{lemma}[Two-slot summary]
At each tree vertex, retain the two least donor records of distinct
original slots, ordered first by donor length, and any two target records
of distinct slots.  These summaries suffice to recover the largest
guaranteed gain of a distinct-slot donor--target pair available at that
vertex.
\end{lemma}
\begin{proof}
Consider the shortest retained donor slot $j$.  If there is a target of a
different slot, one of two distinct target representatives supplies it:
at most one can have slot $j$.  This donor is optimal.  If no such target
exists, every target has slot $j$.  The shortest donor with a different
slot is then the second retained donor; if it does not exist, no eligible
pair exists.  The same argument covers vertices with only one target or
one donor slot.  Donor orientation and occurrence location do not affect
the gain at this fixed depth, so a single representative per donor slot
is sufficient.
\end{proof}

These constant-size summaries merge bottom-up.  The implementation keeps
deterministic tuple orders for witnesses, but does not promise the
historical suffix automaton's tie choice.  Maximum gain, distinct slots,
valid offsets and literal equality are the required invariant properties.

\begin{theorem}[Exact shared-index query]
The shared capped suffix-link index returns a maximum positive guaranteed
decrease over all distinct relator slots, cyclic rotations and donor
signs.  It uses $O(s+L\log(L+1))$ integer and dictionary operations and
$O(L)$ auxiliary storage.  Under the usual expected constant-time
hash-dictionary assumption this is the corresponding expected word-time
bound.
\end{theorem}
\begin{proof}
Every candidate common cyclic segment is represented by two marked points,
by construction and the inverse-window lemma.  Their longest common suffix
ends at their lowest common ancestor in the tree after cap insertion.
The gain is increasing in the common depth when the donor is fixed, so
an optimum is attained at a tree or cap vertex.  The two-slot lemma
recovers an optimal pair at every such vertex.  Taking the largest of
these values therefore gives the global optimum, and each recovered
pair is a literal valid match.

The indexed text has length $2L$.  Building its suffix automaton and
scanning all inverse blocks uses $O(L)$ dictionary operations.  A depth-first
traversal maintains the increasing depths on its current suffix-link path;
binary search locates each of at most $2L$ terminal caps in $O(\log(L+1))$
operations.  Caps on each original edge are sorted.  Their total sorting
cost is $O(L\log(L+1))$, since the sum of their counts is $O(L)$.  Constant-size
summary propagation and candidate evaluation are linear in the capped
tree size.  The original list of $s$ slots is traversed once.  All stored
arrays, transitions, window marks and cap records have linear total size.
\end{proof}

For a node of depth $h$, with donor endpoint $e_d$ and target endpoint
$e_t$, the certificate offsets are
\[
 a_d=(e_d-h+1)\bmod d,\qquad
 a_t=(e_t-h+1)\bmod t.
\]
The resulting record uses the existing version-two relator fields.  No
change to the certificate format or the independent checker is needed.

\subsection{A bounded prelude that improves the actual knot query}

The shared index alone is slower on Gordian than the small set of useful
pairwise scans.  This measured regression motivates the implemented
adaptive query.  As a fixed-size base case, lists with at most four
nonempty relators use the existing pairwise query.  Its cost is $O(s+L)$
in that case, and its historical tie choices are preserved.  The threshold
four is fixed independently of the input; allowing it to grow with the
input would require including it as a parameter in the complexity bound.
This base case avoids paying for both a failed prelude and a shared index
on the two small mirror examples.  For larger lists, every donor of
length $d$ satisfies
\[
 2\min(d,t)-d\le d.
\]
Consequently a donor with $d\le G$, where $G$ is the largest gain already
found, cannot improve the incumbent.  Process donors in decreasing length,
retain the existing target upper bound, and stop once all remaining donor
lengths are at most $G$.

This is a sound search order and pruning rule, but on its own retains the
old worst-case $mL$ scan cost.  The implementation therefore limits the
prelude to
\[
 Q=2\max(1,L)\,\operatorname{bitlength}(L+1)
\]
charged work units.  Sorting, donor construction, suffix-link traversal,
transition creation, clone copying, target-pair consideration and target
scanning consume this allowance.  When the local allowance is exhausted,
the complete shared index is built under the same enclosing resource
budget.  Genuine global cancellation or resource exhaustion propagates;
it is not converted into an invitation to reset the budget.

\begin{corollary}[Adaptive exact query]
The adaptive query has the same maximum-gain guarantee and
$O(s+L\log(L+1))$ dictionary/word-operation bound as the shared query.
It retains $O(L)$ auxiliary storage, including failed prelude construction.
\end{corollary}
\begin{proof}
The fixed-slot base case has linear cost and is exact.  Otherwise, if
the prelude completes, all unvisited donors and targets have valid upper
bounds no larger than the incumbent.  It is exact.  Otherwise the complete
shared query is exact independently of the partial prelude.  Prelude work
is bounded by $Q$ plus one already charged operation batch of $O(L)$; the
shared query then has the preceding theorem's bound.  The one-donor
index is no longer needed when the continuation begins.  No list of
all donor automata or pairwise matches is retained.
\end{proof}

Dictionary-operation bounds should not be confused with unconditional
bounds for CPython's deterministic integer hashing.  A deterministic
balanced-map implementation gives $O(s+L\log(L+1))$ comparison operations
for the shared index alone and a conservative
$O(s+L\log^2(L+1))$ comparison bound for this adaptive scheduling policy.
Comparisons and stored lengths have their ordinary bit costs.  The
presented Python code uses its existing dictionary conventions; it does
not implement a new balanced-map data structure.

On Gordian the maximum guaranteed decrease is $2\cdot3362-3726=2998$.
Only the length-$4260$ and length-$3726$ donors survive the decreasing-length
prelude's final bound.  Four oriented donor automata suffice, versus 24
in the recorded incumbent query.  The selected move and the entire
141-move certificate remain identical on this input.  The timing table
and raw data report both the resulting full-query speedup and the
shared-index-only regression.

\subsection{Scope and remaining research obligations}

This result accelerates an exact query on an explicit presentation.  It
does not establish a small explicit presentation for every input knot,
a bound on grammar growth, a bound on the number of changing relator
moves, or completeness of greedy presentation simplification.  If $L$
becomes exponential in the original crossing number, even the new local
query may take exponential time.  The existing compressed-word pipeline
continues to handle inputs beyond its explicit expansion cap; this
shared explicit index is used only when an explicit query is actually
made.  No false negative is inferred from an absent shortening move.

For a more explicit conditional accounting, suppose the ordinary explicit
producer starts with $r$ generators and its total reduced relator length
never exceeds $B$.  There are at most $r-1$ generator eliminations.  Between
two successive eliminations, every accepted Whitehead or overlap move
strictly decreases an integer length in $[0,B]$, so there are at most $B$
such moves.  Including an unsuccessful final query, the number of overlap
queries is at most $r(B+1)$.  Their total new cost is therefore
\[
 O\!\left(r(B+1)\bigl(s+B\log(B+1)\bigr)\right)
\]
dictionary/word operations.  This charges only the overlap searches, not
Whitehead cuts, generator substitution, replay or diagram preparation.
If $B$ were universally quasi-polynomially bounded, this contribution
would be quasi-polynomial as well.  Establishing such a bound and ensuring
successful recognition remain separate obligations.

A useful next problem is a similarly shared fully compressed cyclic index:
can one remove the factor equal to the number of relators while charging
only grammar size and logarithmic represented lengths?  Another is an
incremental index under a certified relator replacement, with deleted
occurrences handled exactly.  Both require new representation bounds;
copying this explicit tree into a compressed arena would not supply them.
